\documentclass[a4paper,10pt]{article}
\usepackage{lmodern}
\usepackage{cw2026-paper}
\usepackage[italic]{mathastext}

% A short, specific title: identify the behaviour and the investigation.
\newcommand{\PaperTitle}{Sensor footprint and wheel calibration limits on raster-scan image fidelity}

% Replace each placeholder with a seven-digit student number.
\newcommand{\StudentNumbers}{[7-digit student ID]\quad
  [7-digit student ID]\quad [7-digit student ID]\quad [7-digit student ID]}

\begin{document}
\PaperTitleBlock

\begin{multicols*}{2}
\raggedcolumns

\section{Introduction}
A small differential-drive robot carries five reflectance sensors that look at
the floor. It drives up and down the sheet in a coverage pattern~\cite{choset} and
rebuilds a picture from the darkness of each reading. The encoders say where a reading was taken. The
sensors say how dark the paper was. Together they produce the image that the
task is judged on, as Figure~\ref{fig:setup} shows. The digital twin runs the
same controller against a model of the same sheet at one pixel per millimetre.
We use the twin to choose scan settings, so its fidelity matters in practice. A
twin that overstates which features survive recommends settings that
underperform at the bench.

Two explanations are available. The first is optical. Each sensor averages the
paper over a finite spot, so edges must be softened~\cite{richardson,lucy}. The
second is geometric. Encoder odometry drifts because the controller's model of
the wheels is not the robot's~\cite{borenstein,thrun}, so samples are deposited
in the wrong place.

\begin{quote}
\small\textbf{Question.} At a row pitch of 7\,mm and a pulse-width command of 30,
which of three things limits the fidelity of the rebuilt image, the sensor's
footprint, the row pitch we choose, or the controller's belief about wheel
separation?
\end{quote}

We predicted three things. The footprint sets a resolution floor of a few
millimetres. A row pitch below that floor buys nothing. A geometry error large
enough to be measured on the robot costs more fidelity than either.

\PaperFigure[42mm]{figures/fig_setup.pdf}{}{}%
{The printed sheet and the path. Rows are straight and the transitions between
them are curved, so the robot follows one continuous zigzag and never turns on
the spot. The axle starts on the cross facing $+x$. The five sensors sit 28 to
41\,mm ahead of the axle and within 27\,mm of the centre line, and each reading
is deposited at that sensor's position.}
{fig:setup}

\section{Implementation}
The robot is driven by a state machine clocked at 100\,Hz, and
Figure~\ref{fig:flow} shows its logic. A scan starts from rest. The encoders are
zeroed and the robot settles for 300\,ms, so every run begins from the same
pose. Each row is then driven with a 300\,ms ramp. A heading hold with a gain of
60 PWM per radian and a clamp of plus or minus 15 PWM keeps the robot pointing
along the row, and a slower lane controller steers back onto the row line,
\begin{equation}
  \psi^\star = \psi_r + d\,\arctan\!\left(\frac{y_r-y}{L}\right),
  \qquad |\psi^\star-\psi_r|\le 0.45\ \text{rad},
  \label{eq:steer}
\end{equation}
where $\psi_r$ is the row heading, either 0 or $\pi$, $d$ is plus or minus one
for the direction of travel, $y$ is the estimated lateral position and $L$ is a
lookahead of 35\,mm. The lane change is spread over the 50\,mm run-in past each
end of the image, so the robot never turns on the spot. The shape is deliberate.
An in-place turn leaves a heading error that the next row inherits, and on our
robot the scrub was large enough that turning rows arrived between 5 and 32\,mm
apart when 14 was commanded. A curved transition spends the same steering budget
over 50\,mm and never stops, so every row starts from a corrected heading.

\PaperFigure[54mm]{figures/fig_flow.pdf}{}{}%
{State diagram of the controller. Rows alternate direction without turning. The
run-in past each end of the image is where the robot steers onto the next row
line. A 30\,s guard stops the robot and reports a fault if a row makes no
progress.}
{fig:flow}

\subsection*{To recreate one scan}
\begin{enumerate}
\setlength{\itemsep}{0pt}
\item Print the A3 sheet at 100 percent and tape all four corners flat.
\item Put the robot on the start cross, facing the arrow, with the image to its
left.
\item Set the scan to 17 rows, 160\,mm rows, 7\,mm pitch and command 30.
\item Press the button once. The robot waits 1\,s, then scans and stops.
\item Press Download CSV on the dashboard before the ring buffer fills.
\item Load that CSV into the twin and press S. The twin scans the same sheet
with the same controller.
\end{enumerate}

Two values matter most. The 7\,mm pitch is three quarters of the 9\,mm spacing
between the middle sensors, so neighbouring rows interleave. The 50\,mm run-in
gives the lane controller time to settle before the image starts, which is what
keeps the rows parallel without a turn.

Two parameter sets are used. \emph{Nominal} is the template's ideal robot, with
point sensors on a symmetric arc, 90\,mm separation, no deadband and no lag.
\emph{Calibrated} uses our measured separation of 82.5\,mm, a gain that gives the
58\,mm\,s$^{-1}$ we measured at command 30, the measured sensor offsets and a
3\,mm footprint. The controller's beliefs about the wheels are set apart from the
physics, so a real odometry error is reproduced instead of simulated away.

\section{Experimental method}
\subsection*{Variables and conditions}
Table~\ref{tab:variables} lists what was varied, what was measured and what was
held fixed. Each condition was repeated ten times. In the twin a repeat differs
only through the placement error of 1\,mm and the start-heading error of
0.5$^\circ$ applied to the true pose. The sheet, the path, the speed and the
sensor model were held fixed. Battery state was held roughly constant by
charging before each session.

\begin{minipage}{\linewidth}
\captionof{table}{Experiment variables. Every twin condition uses the calibrated
parameter set unless stated. Each row has $n=10$ runs.}
\label{tab:variables}
\small
\renewcommand{\arraystretch}{1.1}
\setlength{\tabcolsep}{4pt}
\begin{tabularx}{\linewidth}{@{}p{17mm}>{\raggedright\arraybackslash}X@{}}
\toprule
\textbf{Type} & \textbf{Description} \\
\midrule
Independent & Row pitch 5, 7, 10, 13.5 and 20\,mm, with rows adjusted so the
scanned height stays 112\,mm. Belief about wheel separation 82.5 to 90\,mm,
which is 0 to 9.1\,\% of the true value. Sensor footprint $\sigma$ from 0 to
5\,mm, run at 3\,mm pitch. U-turn runs at 14\,mm pitch with rotation efficiency
from 0.8 to 1.0. \\
\addlinespace[3pt]
Dependent & Image error $E$ in percent. Coverage $C$ in percent. Pose error $e$
in mm. Edge width $w$ in mm. Row-centre spacing in mm. \\
\addlinespace[3pt]
Controlled & Sheet, $160\times112$\,mm area, 7\,mm pitch unless varied, command
30, 50\,mm run-in, 300\,ms settle, 100\,Hz control, sensor offsets. \\
\bottomrule
\end{tabularx}
\end{minipage}

\subsection*{Metrics}
Let $\hat{I}$ be the rebuilt image and $S$ the sheet, both scaled to $[0,1]$, and
let $P$ be the set of pixels that received at least one sample. Image error,
coverage and pose error are then
\begin{align}
  E&=\frac{100}{|P|}\sum_{p\in P}\left|\hat{I}(p)-S(p)\right|,\label{eq:metrics}\\
  C&=100\,\frac{|P|}{|P_A|},\qquad
  e_k=\left\|\hat{\mathbf{q}}_k-\mathbf{q}_k\right\|,\nonumber
\end{align}
where $|P_A|$ is the nominal scan area in pixels, and $\hat{\mathbf{q}}_k$ and
$\mathbf{q}_k$ are the estimated and true sensor positions of sample $k$. $E$
answers how much of the picture is wrong. $C$ answers how much of it exists.
$e_k$ is reported as a mean $\bar{e}$ and a maximum $e_{\max}$ because a small
bias and a large excursion have different consequences.

Edges are characterised by their width from 10 to 90 percent of the transition.
If the sensor is a Gaussian spot of standard deviation $\sigma$ and the rebuild
deposits each reading with a Gaussian of width $s$, a step edge returns as an
error function and the width is
\begin{equation}
  w = 2.563\sqrt{\sigma^2+s^2} ,
  \label{eq:edge}
\end{equation}
with $s=1$\,mm for our rebuild. This is the line-spread relation used to turn a
measured edge into a footprint estimate~\cite{boreman}. A trial counts as complete if the state machine reaches the end.
No trial failed.

\subsection*{Procedure and analysis}
A trial resets the pose, waits 1\,s, scans and stops. The twin writes one file
of pose, five sensor channels and state per sample, plus an image of the
rebuild. The analysis is a script over those files. The robot contributes
characterisation values instead of scored scans. These are wheel separation from
a calibrated U-turn, 82.5\,mm against the template's 90\,mm, speed from scan
logs, and row spacing measured with a rule after U-turn runs.

\section{Results}
Figure~\ref{fig:path} shows the path the twin drove, Figure~\ref{fig:simreal}
compares the twin with the robot, Figure~\ref{fig:pitch} varies the row pitch and
Figure~\ref{fig:sigma} varies the sensor footprint.

\PaperFigure[36mm]{figures/fig_path.pdf}{}{}%
{The executed path of one twin scan, coloured by row. Rows are straight and each
end is a smooth transition in the run-in, so no row begins with a turn.}
{fig:path}

The template twin scores 15.2 percent with a spread of 4.2 in 123\,s, and the
calibrated twin 20.6 with a spread of 1.0 in 86\,s. They differ in six parameters,
so the gap of 5.4 points is a size and not a cause. A two by two comparison
separates two of them. No footprint with a correct belief gives 13.3.
Adding only the footprint gives 20.9. Adding only the template belief of 90\,mm
gives 21.2. Both give 25.2. The footprint is worth 5.8 points and the belief 6.1,
and the interaction is minus 3.6, so the two errors overlap.

% REAL-SCAN-BEGIN
% REAL-SCAN-END

\PaperFigure[56mm]{figures/fig_simreal.pdf}{}{}%
{The twin against the robot. The upper panel plots pose error against the
odometry model's wheel-separation error, with $n=10$, shading of one standard
deviation and the measured 7.5\,mm marked. The lower panel plots row-centre
spacing for the same 14\,mm command, with the band measured on the robot.}
{fig:simreal}

In the upper panel the mean pose error grows from 1.6 to 5.1\,mm and the maximum
from 2.2 to 10.2\,mm over 0 to 9.1 percent belief error, a fourfold rise in the
maximum, while the image error rises from 21.0 to 25.2 percent. Growth is flat to
3.6 percent and then accelerates. Our measured 7.5\,mm sits at the top of the
curve, which predicts 4 to 5\,mm of mean error and maxima near 10\,mm, larger than
the 7\,mm pitch. In the lower panel the twin gives one line at $14.2\pm0.4$\,mm
with no slip and two bands near 29 and 32\,mm with 5 percent slip, while the robot
spans 5 to 32\,mm. No single efficiency reproduces that spread, and the twin
cannot go below 14\,mm.

\PaperFigure[46mm]{figures/fig_pitch.pdf}{}{}%
{Image error against row pitch, with $n=10$ per point. Blue scores only measured
pixels, red scores the whole target and counts unmeasured area as wrong. Bars are
one standard deviation. The lower panel gives coverage, with scan times.}
{fig:pitch}

Pitch matters only if the missing area counts. On measured pixels the error is
flat, 20.4 to 22.2 percent from 5 to 20\,mm, while coverage falls from 99.8 to
75.6 percent and time from 117 to 34\,s. On the whole target the same runs rise
from 21.3 to 39.8 percent. The metric divides by the pixels that received a
sample, so a quarter of the image is excluded at 20\,mm, and the target rectangle
is also shorter. For a task that must reproduce the sheet, the red series is the
honest one.

\PaperFigure[46mm]{figures/fig_sigma.pdf}{}{}%
{Sensor footprint against edge sharpness at 3\,mm pitch, with $n=10$. The upper
panel gives measured edge width, Equation~\ref{eq:edge} and image error. The band
is our robot's 3\,mm footprint, which is a calibration input. Below, the triangle
at three footprints.}
{fig:sigma}

The edge width grows from 2.6 to 11.7\,mm as the footprint goes from 0 to 5\,mm.
Equation~\ref{eq:edge} predicts 8.1\,mm at 3\,mm and 13.1 at 5\,mm, against 7.7
and 11.7 measured, so the spot is not perfectly Gaussian. The image error rises
from 14.8 to 25.0 percent.

The U-turn runs show the metric failing. From efficiency 1.0 to 0.8 the mean pose
error grows from 2 to 247\,mm, yet the image error peaks at 39.3 percent at 0.98
and falls to 30.8. Each channel is rescaled to its own range and the error counts
only covered cells, so a scan that sees little is stretched and judged on what it
reached.

\section{Discussion and conclusion}
\subsection*{Interpretation}
On measured pixels the footprint costs 5.8 points and the belief 6.1, and pitch
costs nothing. On the whole target pitch costs 18.5. The ranking depends on the
metric, and only the whole-target one matches the task. If the robot must deliver
a picture, the missing quarter is error, and pitch is the largest effect.

The prediction was half right. The footprint sets a floor of about 8\,mm, but
geometry does not cost more than pitch. The twin is optimistic in a narrower sense
than we claimed. An ideal sensor scores 15.2 against 20.6 for the calibrated twin,
even though the template also has the wrong belief, so that score says what an
ideal sensor can see and not what the bench will do.

Two alternatives remain. The flat curve could be saturation, since detail between
1 and 3\,mm is already lost. The robot's spread could be placement error, which a
rule cannot separate.

The trade-off is explicit. Seven millimetres gives 99.8 percent coverage in 86\,s.
Ten millimetres gives 92.5 percent in 60\,s, with the same error, 21.0 against 21.5,
and a lower pose error, 1.2 against 2.1. If 92.5 percent is enough, 10 wins. We
chose 7 because the task is to reproduce the sheet.

\subsection*{Limitations and next experiment}
The limitation that most affects confidence is that the real side is a parameter
set, not a scored scan. No rebuild from the robot is scored anywhere here, so the
5.4-point difference is model to model. It shows how far the model moves when the
template is replaced. It does not show how far either is from the bench.

Three limits qualify the numbers. The repeats are rigid re-registrations of one
deterministic scan, since only the start pose is randomised, so the bars measure
sensitivity to placement. Several calibrated parameters, including the gain and
the footprint, are estimates. And the error metric averages over a population that
changes with the condition.

The next experiment follows directly. Scan the sheet three times with the robot and
score each log with the same rebuild. That turns the model-to-model gap into a
model-to-data gap. A black and white edge at three speeds would then separate the
footprint from the delay, which is currently zero.

\subsection*{Conclusion}
The reconstruction is limited by three things at different scales. On the pixels
sampled, the spot costs 5.8 points and the wheel belief 6.1. Across the whole
target, pitch costs 18.5 because it decides how much is sampled. The twin's default
metric hides the second, so a twin used to choose settings should score the full
target. The robot is harder than its twin in a way the twin cannot show. Its row
spacing spreads from 5 to 32\,mm where the model gives a line.

\newpage
\section*{Evidence and sources}
Twin runs are in \texttt{all\_runs.csv}, with 300 rows and one row per twin run,
written by \texttt{run\_experiments.py}. It holds every condition used here,
including the full U-turn sweep from 0.8 to 1.0 and the interaction cell with no
footprint and the template belief. Per-run samples are in
\texttt{exp/<condition>/runs/} and the rebuilds are in the matching
\texttt{.png} files. The figures are produced by \texttt{analyse.py}, which reads
only those files. The twin source is \texttt{DigitalTwin/}, where
\texttt{ModelParameters.pde} holds the nominal and the calibrated sets. The robot
firmware is in \texttt{Controller.cpp}, \texttt{ScanLog.cpp} and
\texttt{dashboard/dashboard.html}. The headless driver used to batch the twin is
\texttt{merge.py} plus a batch tab, and it does not change the behaviour of the
sketch. Measured robot values are recorded in \texttt{ModelParameters.pde} and in
the repository README. The twin is written in Processing~\cite{processing}.

\renewcommand{\refname}{}
\begin{thebibliography}{9}
\small
\bibitem{richardson} W.~H. Richardson, ``Bayesian-based iterative method of image
restoration,'' \emph{J. Opt. Soc. Am.} \textbf{62}(1), 55--59, 1972.
\bibitem{lucy} L.~B. Lucy, ``An iterative technique for the rectification of
observed distributions,'' \emph{Astron. J.} \textbf{79}(6), 745--754, 1974.
\bibitem{borenstein} J.~Borenstein and L.~Feng, ``Measurement and correction of
systematic odometry errors in mobile robots,'' \emph{IEEE Trans. Robotics
Autom.} \textbf{12}(6), 869--880, 1996.
\bibitem{thrun} S.~Thrun, W.~Burgard and D.~Fox, \emph{Probabilistic Robotics},
MIT Press, 2005, ch.~5.
\bibitem{choset} H.~Choset, ``Coverage for robotics, a survey of recent
results,'' \emph{Ann. Math. Artif. Intell.} \textbf{31}, 113--126, 2001.
\bibitem{boreman} G.~D. Boreman, \emph{Modulation Transfer Function in Optical
and Electro-Optical Systems}, SPIE Press, 2001, ch.~3.
\bibitem{tao} F.~Tao, H.~Zhang, A.~Liu and A.~Y.~C. Nee, ``Digital twin in
industry, state of the art,'' \emph{IEEE Trans. Ind. Informat.}
\textbf{15}(4), 2405--2415, 2019.
\bibitem{processing} Processing Foundation, \emph{Processing 4 reference},
processing.org, accessed 2026.
\end{thebibliography}

\end{multicols*}
\end{document}
